\chapter{Analiza cerintelor}

\section{Parti interesate si roluri}

Trainee este construit in jurul a trei roluri principale:
\begin{itemize}[leftmargin=1.2cm]
    \item \textbf{Client}: cauta antrenori, consulta profiluri, urmareste programari si foloseste fluxul de check-in;
    \item \textbf{Antrenor}: isi administreaza profilul, disponibilitatea si alocarea clientilor pe intervale;
    \item \textbf{Administrator}: gestioneaza incidentele si operatiunile de moderare.
\end{itemize}

Pe langa aceste roluri, exista parti interesate secundare: echipa tehnica, zona de suport si furnizorii de servicii externe (billing, email, stocare).

\section{Cerintele functionale}

\subsection{Autentificare si cont}

\begin{enumerate}[leftmargin=1.2cm]
    \item Sistemul permite inregistrare cu email si parola.
    \item Sistemul permite autentificare si emitere token valid.
    \item Sistemul permite resetare parola prin email.
    \item Sistemul permite verificarea adresei de email.
\end{enumerate}

\subsection{Descoperire antrenori}

\begin{enumerate}[leftmargin=1.2cm]
    \item Clientul poate cauta antrenori dupa text, localizare, tarif si experienta.
    \item Rezultatele pot fi sortate dupa criterii functionale, inclusiv distanta.
    \item Sistemul afiseaza profil complet pentru fiecare antrenor.
    \item Sistemul limiteaza inflatarea artificiala a vizualizarilor de profil.
\end{enumerate}

\subsection{Scheduling}

\begin{enumerate}[leftmargin=1.2cm]
    \item Antrenorul poate defini intervale de lucru pe zile.
    \item Sistemul poate genera sloturi pe perioade calendaristice.
    \item Antrenorul poate aloca clienti pe sloturi disponibile.
    \item Clientul poate folosi coduri de check-in in fluxul dedicat.
\end{enumerate}

\subsection{Billing}

\begin{enumerate}[leftmargin=1.2cm]
    \item Sistemul gestioneaza entitlement-ul de subscriere pentru antrenori.
    \item Sistemul proceseaza webhook-uri in mod idempotent.
    \item Sistemul permite validarea achizitiilor in-app pe platforme mobile.
\end{enumerate}

\subsection{Suport si administrare}

\begin{enumerate}[leftmargin=1.2cm]
    \item Utilizatorii pot raporta incidente.
    \item Administratorii pot filtra si actualiza statusul incidentelor.
    \item Sistemul pastreaza date suficiente pentru audit operational.
\end{enumerate}

\section{Cerintele nefunctionale}

\begin{table}[H]
\centering
\caption{Cerintele nefunctionale principale}
\begin{tabular}{p{0.22\textwidth} p{0.68\textwidth}}
\toprule
\textbf{Categorie} & \textbf{Criteriu} \\
\midrule
Performanta & Endpoint-urile de cautare trebuie sa raspunda stabil pentru trafic uzual; paginarea si limitarea rezultatului sunt obligatorii. \\
Securitate & Endpoint-urile publice necesita rate limiting; payload-urile sunt validate strict; secretele nu se expun in client. \\
Disponibilitate & Aplicatia trebuie sa degradeze controlat cand servicii externe devin indisponibile partial. \\
Scalabilitate & Arhitectura trebuie sa permita extinderea pe mai multe instante si inlocuirea store-ului de rate limiting. \\
Mentenabilitate & Modulele trebuie separate pe feature-uri, cu tipuri explicite si configuratie centralizata. \\
Observabilitate & Sistemul trebuie sa poata integra health checks, metrici de latenta si monitorizare trafic. \\
Conformitate & Produsul trebuie sa includa documente legale si reguli de protectie a datelor. \\
\bottomrule
\end{tabular}
\end{table}

\section{Constrangeri de proiect}

Cea mai importanta constrangere a fost pastrarea functionalitatii pe durata evolutiei tehnice. Cu alte cuvinte, masurile de hardening nu trebuiau sa rupa fluxurile deja folosite in aplicatie.

Alte constrangeri relevante:
\begin{itemize}[leftmargin=1.2cm]
    \item dependenta de servicii externe cu reguli proprii (Stripe, RevenueCat, AWS, SMTP);
    \item interval de timp limitat pentru implementare si documentare;
    \item necesitatea alinierii la regulamentele academice de finalizare studii.
\end{itemize}

\section{Scenarii de utilizare reprezentative}

\subsection{Scenariul S1: Cautare antrenor}

Clientul acceseaza ecranul de cautare, aplica filtre dupa locatie si tarif, apoi intra pe profilul unui antrenor. Backend-ul valideaza query-ul, executa filtrarea si inregistreaza evenimentul de vizualizare conform regulilor anti-abuz.

\subsection{Scenariul S2: Programare asistata de cod}

Antrenorul genereaza sloturi pentru perioada curenta, selecteaza un interval liber si asociaza clientul prin cod de check-in. Sistemul verifica validitatea codului, marcheaza asocierea si actualizeaza statusul slotului.

\subsection{Scenariul S3: Raportare incident}

Utilizatorul transmite un incident din aplicatie, iar administratorul il preia, il analizeaza si actualizeaza statusul pana la rezolvare.

\section{Model de cazuri de utilizare}

\begin{figure}[H]
\centering
\fbox{\parbox{0.88\textwidth}{
\textbf{Placeholder diagrama Use Case} \\
Diagrama finala va include actorii Client, Antrenor si Administrator, impreuna cu cazurile-cheie: autentificare, cautare antrenori, programari, billing, raportare si administrare incidente.
}}
\caption{Model conceptual de cazuri de utilizare}
\end{figure}

\section{Criterii de acceptanta}

Criteriile de acceptanta definesc pragul minim pentru validarea implementarii:
\begin{enumerate}[leftmargin=1.2cm]
    \item Fluxurile de autentificare functioneaza end-to-end.
    \item Cautarea antrenorilor produce rezultate coerente pentru filtre valide.
    \item Scheduling-ul permite creare si alocare de sloturi fara erori de integritate.
    \item Endpoint-urile publice aplica limitare de trafic si validare de payload.
    \item Billing-ul sincronizeaza entitlement-ul pe baza evenimentelor externe.
    \item Modulul de issue management acopera ciclul complet de raportare si rezolvare.
\end{enumerate}

Prin formalizarea acestor cerinte, capitolul urmator poate detalia arhitectura tehnica aleasa pentru a le sustine.
